\subsection{Problem formulation}
For economy $c$ and year $t$ let $\mathbf{x}_{c,t}\in\mathbb{R}^{p}$ be the socio-economic, fiscal, demographic and governance features, and let $\mathbf{y}_{c,t}=(y_1,\dots,y_{K})$ be the percentage shares of total general-government expenditure assigned to $K=17$ components (16 COFOG functions plus the residual; $K=10$ for panel A10), so that $y_k\ge 0$ and $\sum_k y_k=100$. The task is to learn $\hat{\mathbf{y}}_{c,t}=g(\mathbf{x}_{c,t},\mathbf{h}_{c,t})$ where $\mathbf{h}_{c,t}$ is optional history (last year's shares). The output is a descriptive reference for a human analyst, not a decision.

\textbf{Information set (nowcast, not forecast).} The features $\mathbf{x}_{c,t}$ are the \emph{same-year} values of GDP growth, inflation, debt and so on. These are realised values that are published with a lag and are not known when a budget for year $t$ is drafted; moreover, official COFOG for $t-1$ is itself published with a lag of 12--24 months. The annual regime below is therefore a \emph{nowcast} conditional on realised covariates, and persistence (which uses only $\mathbf{y}_{c,t-1}$) is evaluated on a less favourable information set than the learners. A strict forecast with lagged covariates was not run.

\figArch

\subsection{Prediction layer}
Two deployment regimes are distinguished because they correspond to different situations.
\paragraph{Regime~1: annual update with history.} The analyst knows last year's allocation $\mathbf{y}_{c,t-1}$. We learn the \emph{change}
\begin{equation}
\hat{\mathbf{y}}_{c,t}=\mathbf{y}_{c,t-1}+f\big(\mathbf{x}_{c,t},\mathbf{y}_{c,t-1}\big),
\end{equation}
so that persistence is the null model and the learner must earn any improvement.
\paragraph{Regime~2: cold start (no history).} For an economy without COFOG history, $\hat{\mathbf{y}}_{c,t}=f(\mathbf{x}_{c,t})$ is learned from \emph{other} economies; the benchmark is the mean allocation of the same World Bank income group (four classes).
\paragraph{Learners.} One regressor per component for XGBoost~\cite{chen2016xgboost}, LightGBM~\cite{ke2017}, CatBoost~\cite{prokhorenkova2018} and Random Forest~\cite{breiman2001}, and a ridge regression as a linear reference (hyper-parameters in Table~\ref{tab:hyper}; targets standardised before fitting). Tree ensembles are chosen because they are strong on medium-sized tabular data~\cite{grinsztajn2022}; our panels are smaller than those in that benchmark, so this is a prior, not a guarantee. The \emph{hybrid} is the equal-weight mean of XGBoost and CatBoost predictions, $\hat{\mathbf{y}}^{\text{hyb}}=\tfrac12(\hat{\mathbf{y}}^{\text{xgb}}+\hat{\mathbf{y}}^{\text{cat}})$; it is evaluated against its parts rather than assumed superior. Its choice was fixed before the test window was scored, but five learners and several variants were evaluated, so a best-of-several selection effect cannot be excluded.

\subsection{Further predictors evaluated}
\textbf{Log-ratio targets and stacking (architecture v2.1).} (i) \emph{History features:} lag-1 and lag-2 shares, their difference (trend) and the expanding mean of the country's \emph{past} shares. (ii) \emph{Compositional targets:} the centred log-ratio $\mathrm{clr}(\mathbf y)=\log(\mathbf y+\varepsilon)-\overline{\log(\mathbf y+\varepsilon)}$ with $\varepsilon=0.05$~\cite{aitchison1986,egozcue2003}, inverted by a softmax so that predictions are positive and sum to 100. (iii) \emph{Stacking}~\cite{wolpert1992}: for every component $k$ the final prediction is $\sum_m w_{m,k}\hat y^{(m)}_k$ with non-negative weights found by non-negative least squares and normalised to sum to one. In the annual mode the members are persistence, the past country mean and XGBoost/CatBoost learners (raw and log-ratio); weights are fitted on a validation window (train $\le$2016, validate 2017--19), the base models are refit on data up to 2019 and the test window is used once. In the cold-start mode the members are the peer mean, ridge, and XGBoost/CatBoost on log-ratio targets; weights come from a nested inner grouped cross-validation inside every outer training fold, so no test economy influences its weights.
\textbf{Nearest-economies predictor.} Features are median-filled and standardised with training economies only; each training economy is represented by the median of its feature vector over years; the query is compared with every training economy by Euclidean distance and the prediction is the inverse-distance-weighted average of the mean compositions of the $k$ nearest economies ($k=8$ in the comparisons; chosen from $\{3,5,8,12\}$ on out-of-fold data). Its neighbours are shown to the user, which makes it directly explainable.
\textbf{Dirichlet regression and neural networks.} An L2-penalised Dirichlet regression with concentration parameters $\alpha_k=\exp(\mathbf w_k^{\top}\mathbf x+b_k)$ on standardised features (a compositional model), a multilayer perceptron (five-seed ensemble, early stopping on an inner validation split of training economies) and a small feature-tokenising transformer~\cite{gorishniy2021} serve as comparators; both neural models predict softmax shares.
\textbf{Intervals (split conformal).} Absolute residuals per component are computed out-of-fold on the \emph{training} economies (inner grouped folds); their finite-sample $(1-\alpha)$ quantile~\cite{lei2018,angelopoulos2023}, either global or by income group (Mondrian variant), gives a symmetric band around each predicted share. Because out-of-fold residuals of grouped folds are not exchangeable with those of a genuinely new economy~\cite{barber2023}, nominal coverage is not guaranteed and is measured.
\textbf{Applicability score.} The mean standardised Euclidean distance of an economy's feature vector to its five nearest training economies; it is a novelty measure, not a probability.
\textbf{Econometric baselines.} A random walk with country drift, a pooled AR(1) on log-ratio shares with macro covariates, and the same with country fixed effects.

\subsection{Allocation engine (evaluated component)}
Raw predictions need not lie on the simplex or respect user-set limits. The engine first clips negatives and renormalises, then projects onto a bounded simplex:
\begin{equation}
\mathbf{a}^\star=\arg\min_{\mathbf{a}}\ \tfrac12\|\mathbf{a}-\hat{\mathbf{y}}\|_2^2\quad\text{s.t.}\quad \mathbf{1}^{\!\top}\mathbf{a}=100,\ \ \boldsymbol{\ell}\le\mathbf{a}\le\mathbf{u}.
\end{equation}
This is a standard Euclidean projection with solution $\mathbf{a}^\star=\operatorname{clip}(\hat{\mathbf{y}}+\lambda\mathbf{1},\boldsymbol{\ell},\mathbf{u})$ and scalar $\lambda$ found by bisection; we verify it against SLSQP (maximum discrepancy \NSlsqpDiff{} pp), which is a unit test rather than a result. The bounds $\boldsymbol{\ell},\mathbf{u}$ are set from training-data quantiles or minimum/maximum and are where statutory minima would be entered; no statutory constraints are available in our data, so we only study how bound width affects error (Section~\ref{sec:engine}).

\subsection{Explainability layer}
\textbf{SHAP.} Exact TreeSHAP~\cite{lundberg2017,lundberg2020} attributions $\phi_{j}$ are computed for every component model of the XGBoost member; global importance is $\overline{|\phi_j|}$. KernelSHAP (30 held-out rows per division, a 30-centre k-means background of the training rows, 200 samples) is used to explain the stacked blend directly. \textbf{LIME.} Reference tabular LIME~\cite{ribeiro2016} with continuous (non-discretised) features, 1\,500 perturbations and a local ridge surrogate; the signed local contribution of feature $j$ is $w_j(x_j-\bar{x}_j)/s_j$ so that it is comparable with $\phi_j$. Because post-hoc explanations can be unstable or manipulated~\cite{alvarezmelis2018,slack2020} and are not causal, we \emph{measure} them: (i) \emph{agreement} (Spearman $\rho$ between $|\phi|$ and $|$LIME$|$, top-5 Jaccard overlap, sign agreement on the top-5 features); (ii) \emph{stability} (Spearman $\rho$ between two LIME runs with different seeds); (iii) \emph{local fidelity} (LIME surrogate $R^2$); (iv) a \emph{deletion test}: the mean absolute prediction change when the three most important features (by SHAP) are replaced by the training mean, relative to three random features; and (v) a \emph{synthetic ground-truth test} with known drivers. The deletion test uses mean replacement, which creates off-manifold inputs, and any reasonable importance ranking would be expected to pass it; we did not compare with permutation or gain importance.

\subsection{Evaluation protocol}
\textbf{Metrics.} The primary metric is the \emph{total-variation distance} $\mathrm{TVD}=\tfrac12\sum_k|\hat y_k-y_k|$ in percentage points: the share of the budget that would have to be moved to turn the prediction into the observed allocation. It is interpretable to practitioners, scale-free across countries, and a natural reallocation measure; it is \emph{not} a log-ratio (Aitchison) distance, which would weight small functions more and which we do not report. TVD is dominated by the large functions (social protection, health, general public services). With $K$ components, MAE equals $2\,\mathrm{TVD}/K$ exactly, so MAE is not reported separately; RMSE and, for the single-target WDI series, $R^2$ are given where informative. Persistence is the naive benchmark recommended for forecast comparison~\cite{hyndman2006}. \textbf{Splits.} (a) \emph{Temporal}: train 2000--2019, test 2020--2024, which contains the COVID-19 shock and the 2022 energy crisis (block-wise validation respects temporal dependence~\cite{roberts2017,bergmeir2018}); it is a single origin, and persistence is about 60\% worse there than in the 2017--19 validation window (3.62 vs.\ 2.26\,pp); (b) \emph{leave-country-out}: six folds with all years of a country held out, repeated with two fold seeds (five for the robustness study); (c) \emph{leave-region-out} on the WDI panel (Panel B uses one fold seed). \textbf{Uncertainty.} 95\,\% confidence intervals are cluster-bootstrap intervals that resample \emph{countries} (2\,000 draws; 10\,000 in the Holm study) because years within a country are not independent~\cite{cameron2008}; with few clusters (36 in A16; three low-income economies in A10) they are approximate. \textbf{Multiplicity.} The family of eight cold-start comparisons (four per panel) is corrected with the Holm step-down procedure~\cite{holm1979}; the temporal comparison is reported with an uncorrected interval and is not claimed to be significant. \textbf{Paired gain.} Gains are paired differences of row-level TVD, averaged over fold seeds where stated; ``fold-seed~0'' marks values computed on the first fold seed only. \textbf{Baselines.} Persistence, country mean, income-group peer mean (four classes), global mean and ridge; no region-by-income mean or GDP-per-capita regression baseline was run. \textbf{Leakage control.} No feature is a transformation of a target; the audit of the v1 data (Section~\ref{sec:audit}) follows the leakage taxonomy of~\cite{kaufman2012}. Remaining leakage-type issues (transductive imputation, features that are components of spending) are listed in Section~\ref{sec:disc}.

\paragraph{Implementation details.} All experiments ran on a 12-core CPU workstation without GPU, Python 3.13 with XGBoost 3.2, LightGBM 4.6, CatBoost 1.2, scikit-learn 1.8, SHAP 0.52, LIME, PyTorch 2.9 (CPU), pandas 2.2, NumPy 2.2 and SciPy 1.16. Fixed seeds (42; five fold seeds for the robustness study) make every table reproducible; hyper-parameters are listed in the supplementary tables and were not tuned on test folds (Optuna tuning, Section~\ref{sec:ext}, used nested grouped cross-validation inside training economies).
